
    %%%%%%%%%%%%%%%%%%%%%%%%%%%%%%%%%%%%%%%%%%%%%%%%%%%%%%%%%%%%%%%%%%%%%%%%%%%%%%%%
    %%% Classe do documento
    \documentclass[12pt, addpoints]{exam}
    \UseRawInputEncoding

    %%%%%%%%%%%%%%%%%%%%%%%%%%%%%%%%%%%%%%%%%%%%%%%%%%%%%%%%%%%%%%%%%%%%%%%%%%%%%%%%
    % preambulo.tex
    %
    % Arquivo de configuração de packages para uso o LaTeX, conforme minhas
    % preferências e modelos pessoais.
    %
    % Para maiores informações, visite:
    %    https://github.com/abrantesasf/latex
    %
    % NÃO ALTERE SE NÃO SOUBER O QUE ESTÁ FAZENDO!
    %%%%%%%%%%%%%%%%%%%%%%%%%%%%%%%%%%%%%%%%%%%%%%%%%%%%%%%%%%%%%%%%%%%%%%%%%%%%%%%%


    %%%%%%%%%%%%%%%%%%%%%%%%%%%%%%%%%%%%%%%%%%%%%%%%%%%%%%%%%%%%%%%%%%%%%%%%%%%%%%%%
    %%% Estruturas de controle:
    \input{static/utils/controles}

    %%%%%%%%%%%%%%%%%%%%%%%%%%%%%%%%%%%%%%%%%%%%%%%%%%%%%%%%%%%%%%%%%%%%%%%%%%%%%%%%
    %%% Compilação condicional em PDF:
    \input{static/utils/pdf}

    %%%%%%%%%%%%%%%%%%%%%%%%%%%%%%%%%%%%%%%%%%%%%%%%%%%%%%%%%%%%%%%%%%%%%%%%%%%%%%%%
    %%% Configurações de layout da página:
    \input{static/utils/layout}

    %%%%%%%%%%%%%%%%%%%%%%%%%%%%%%%%%%%%%%%%%%%%%%%%%%%%%%%%%%%%%%%%%%%%%%%%%%%%%%%%
    %%% Configurações para autores:
    \input{static/utils/autores}

    %%%%%%%%%%%%%%%%%%%%%%%%%%%%%%%%%%%%%%%%%%%%%%%%%%%%%%%%%%%%%%%%%%%%%%%%%%%%%%%%
    %%% Cabeçalho e rodapé:
    %%% Atenção! É MUITO DIFÍCIL ajustar apropriadamente os running headers
    %%% e running footers da primeira vez. Você pode confiar no LaTeX e deixar que
    %%% ele faça o ajuste automaticamente ou pode quebrar a cabeça e
    %%% tentar melhorar algo que o LaTeX já faz muito bem, mesmo que não fique
    %%% exatamente do jeito que você quer. O que você prefere? Se quiser ajustar
    %%% manualmente, descomente a linha abaixo e faça as configurações no arquivo
    %%% "utils/headings.tex".
    %\input{static/utils/headings}

    %%%%%%%%%%%%%%%%%%%%%%%%%%%%%%%%%%%%%%%%%%%%%%%%%%%%%%%%%%%%%%%%%%%%%%%%%%%%%%%%
    %%% Configuração para as fontes do documento:
    %%% Se você quiser usar fontes próprias ou do sistema, precisará ajustar as
    %%% configurações no arquivo "utils/fontes.tex".
    %%% ATENÇÃO: por padrão eu utilizo XeLaTeX com fontes proprietárias projetadas
    %%% por Matthew Butterick, adquiridas em https://mbtype.com, que não podem ser
    %%% distribuídas devido à licença de utilização. Se você usar XeLaTeX, você
    %%% DEVERÁ OBRIGATORIAMENTE ajustar as fontes no arquivo "utils/fontes.tex".
    \input{static/utils/fontes}

    %%%%%%%%%%%%%%%%%%%%%%%%%%%%%%%%%%%%%%%%%%%%%%%%%%%%%%%%%%%%%%%%%%%%%%%%%%%%%%%%
    %%% Sumário:
    \input{static/utils/toc}

    %%%%%%%%%%%%%%%%%%%%%%%%%%%%%%%%%%%%%%%%%%%%%%%%%%%%%%%%%%%%%%%%%%%%%%%%%%%%%%%%
    %%% Matemática:
    \input{static/utils/matematica}

    %%%%%%%%%%%%%%%%%%%%%%%%%%%%%%%%%%%%%%%%%%%%%%%%%%%%%%%%%%%%%%%%%%%%%%%%%%%%%%%%
    %%% Notas de rodapé:
    \input{static/utils/footnotes}

    %%%%%%%%%%%%%%%%%%%%%%%%%%%%%%%%%%%%%%%%%%%%%%%%%%%%%%%%%%%%%%%%%%%%%%%%%%%%%%%%
    %%% Referências cruzadas, links, citações:
    \input{static/utils/crossrefs}

    %%%%%%%%%%%%%%%%%%%%%%%%%%%%%%%%%%%%%%%%%%%%%%%%%%%%%%%%%%%%%%%%%%%%%%%%%%%%%%%%
    %%% Configurações para os hiperlinks em PDF:
    \input{static/utils/pdfconfig}

    %%%%%%%%%%%%%%%%%%%%%%%%%%%%%%%%%%%%%%%%%%%%%%%%%%%%%%%%%%%%%%%%%%%%%%%%%%%%%%%%
    %%% Glossário e índice remissivo:
    \input{static/utils/glossindex}

    %%%%%%%%%%%%%%%%%%%%%%%%%%%%%%%%%%%%%%%%%%%%%%%%%%%%%%%%%%%%%%%%%%%%%%%%%%%%%%%%
    %%% Referências bibliográficas:
    \input{static/utils/bibliografia}

    %%%%%%%%%%%%%%%%%%%%%%%%%%%%%%%%%%%%%%%%%%%%%%%%%%%%%%%%%%%%%%%%%%%%%%%%%%%%%%%%
    %%% Cores:
    \input{static/utils/cores}

    %%%%%%%%%%%%%%%%%%%%%%%%%%%%%%%%%%%%%%%%%%%%%%%%%%%%%%%%%%%%%%%%%%%%%%%%%%%%%%%%
    %%% Computação:
    \input{static/utils/computacao}

    %%%%%%%%%%%%%%%%%%%%%%%%%%%%%%%%%%%%%%%%%%%%%%%%%%%%%%%%%%%%%%%%%%%%%%%%%%%%%%%%
    %%% Gráficos:
    \input{static/utils/graficos}

    %%%%%%%%%%%%%%%%%%%%%%%%%%%%%%%%%%%%%%%%%%%%%%%%%%%%%%%%%%%%%%%%%%%%%%%%%%%%%%%%
    %%% Ícones:
    \input{static/utils/icones}

    %%%%%%%%%%%%%%%%%%%%%%%%%%%%%%%%%%%%%%%%%%%%%%%%%%%%%%%%%%%%%%%%%%%%%%%%%%%%%%%%
    %%% Blocos:
    \input{static/utils/blocos}

    %%%%%%%%%%%%%%%%%%%%%%%%%%%%%%%%%%%%%%%%%%%%%%%%%%%%%%%%%%%%%%%%%%%%%%%%%%%%%%%%
    %%% Boxes coloridos:
    \input{static/utils/colorbox}

    %%%%%%%%%%%%%%%%%%%%%%%%%%%%%%%%%%%%%%%%%%%%%%%%%%%%%%%%%%%%%%%%%%%%%%%%%%%%%%%%
    %%% Tabelas:
    \input{static/utils/tabelas}

    %%%%%%%%%%%%%%%%%%%%%%%%%%%%%%%%%%%%%%%%%%%%%%%%%%%%%%%%%%%%%%%%%%%%%%%%%%%%%%%%
    %%% Ambientes de listas:
    \input{static/utils/listas}

    %%%%%%%%%%%%%%%%%%%%%%%%%%%%%%%%%%%%%%%%%%%%%%%%%%%%%%%%%%%%%%%%%%%%%%%%%%%%%%%%
    %%% Ambientes floats e similares:
    \input{static/utils/floats}

    %%%%%%%%%%%%%%%%%%%%%%%%%%%%%%%%%%%%%%%%%%%%%%%%%%%%%%%%%%%%%%%%%%%%%%%%%%%%%%%%
    %%% Ajustes para a classe EXAM:
    \input{static/utils/exam}

    %%%%%%%%%%%%%%%%%%%%%%%%%%%%%%%%%%%%%%%%%%%%%%%%%%%%%%%%%%%%%%%%%%%%%%%%%%%%%%%%
    %%% Ajustes para textos:
    \input{static/utils/texto}

    %%%%%%%%%%%%%%%%%%%%%%%%%%%%%%%%%%%%%%%%%%%%%%%%%%%%%%%%%%%%%%%%%%%%%%%%%%%%%%%%
    %%% Ambientes, aliases, comandos e customizações em geral para serem
    %%% utilizadas nos documentos. Defina aqui qualquer customização específica
    %%% para seus documentos!
    \input{static/utils/customizacoes}

    %%%%%%%%%%%%%%%%%%%%%%%%%%%%%%%%%%%%%%%%%%%%%%%%%%%%%%%%%%%%%%%%%%%%%%%%%%%%%%%%
    %%% Ajuste do layout, espaçamento de linhas e etc.:
    \geometry{a4paper, portrait, left=2.0cm, right=2.0cm, top=2.0cm, bottom=2.0cm}
    \singlespacing

    %%%%%%%%%%%%%%%%%%%%%%%%%%%%%%%%%%%%%%%%%%%%%%%%%%%%%%%%%%%%%%%%%%%%%%%%%%%%%%%%
    %%% Configurações de cabeçalho e rodapé (não use fancyhdr pois ocorre conflito):
    \pagestyle{headandfoot}
    \runningheadrule
    \firstpageheader{}{}{}
    \runningheader{banco de dados}{}{Avaliação}
    \firstpagefooter{}{}{Página \thepage\ de \numpages}
    \runningfooter{}{}{Página \thepage\ de \numpages}
    \runningfootrule

    %%%%%%%%%%%%%%%%%%%%%%%%%%%%%%%%%%%%%%%%%%%%%%%%%%%%%%%%%%%%%%%%%%%%%%%%%%%%%%%%
    %%% Configurações para as propriedades do PDF:
    \hypersetup{
    hidelinks,           % Comente para web, descomente para impressão
    colorlinks=true,      % True para web, False para impressão
    pdftitle=BD: Avaliação,
    pdfauthor=Matheus Endlich Silveira,
    pdfsubject=banco de dados,
    pdfkeywords=['Sem tag'],
    pdfinfo={
        CreationDate={}, % Ex.: D:AAAAMMDDHH24MISS
        ModDate={}       % Ex.: D:AAAAMMDDHH24MISS
    }
    }
    \input{static/utils/pdfconfig}

    %%%%%%%%%%%%%%%%%%%%%%%%%%%%%%%%%%%%%%%%%%%%%%%%%%%%%%%%%%%%%%%%%%%%%%%%%%%%%%%%
    %%% Começa o documento
    \begin{document}

    %%%%%%%%%%%%%%%%%%%%%%%%%%%%%%%%%%%%%%%%%%%%%%%%%%%%%%%%%%%%%%%%%%%%%%%%%%%%%%%%
    %%% Folha de instruções padronizadas da UVV para a prova
    \pdfbookmark[1]{Instruções}{instrucoes}
    \begin{figure}[H]
        \begin{center}
            \includegraphics[scale=0.3]{{static/imgs/logo-uvv.png}}
        \end{center}	
    \end{figure}

    \vspace{-1cm}

    \begin{table}[H]
        \centering
        \begin{tabular}{|llll|}
            \hline
            \multicolumn{2}{|l|}{\textbf{Disciplina:} banco de dados} & 
            \multicolumn{1}{l|}{\multirow{3}{*}{\textbf{\makecell{Nota:\\ \,\\ \,}}}} & 
            \multirow{3}{*}{\textbf{\makecell{Coordenador:\\ \,\\ \,}}} \\ 
            \cline{1-2}
            \multicolumn{2}{|l|}{\textbf{Professor:} Matheus Endlich Silveira \hspace{4.72cm} } & 
            \multicolumn{1}{l|}{} & \\ 
            \cline{1-2}
            \multicolumn{2}{|l|}{\textbf{Aluno:}} & 
            \multicolumn{1}{l|}{} & \\ 
            \hline
            \multicolumn{1}{|l|}{\textbf{Turma:} \hspace{6.3cm} } & 
            \multicolumn{1}{l|}{\textbf{Semestre:}} & 
            \multicolumn{2}{l|}{\textbf{Valor:} 10 pontos} \\ 
            \hline
            \multicolumn{1}{|l|}{\textbf{Data:}} & 
            \multicolumn{3}{l|}{\textbf{Avaliação:}} \\ 
            \hline
        \end{tabular}
    \end{table}

    \begin{center}
        \fbox{\setlength\fboxsep{0.3cm} \fbox{\parbox{15.4cm}{
            \begin{center}
                \textbf{INSTRUÇÕES PARA A REALIZAÇÃO DESTA AVALIAÇÃO:}
            \end{center}
            \begin{itemize}[leftmargin=*]
                \item Esta é a primeira avaliação da disciplina banco de dados, 
                    e engloba o conteúdo referente Sem tag.
                \item Esta avaliação contém 12 questões objetivas, \textbf{todas obrigatórias}.
                \item \textbf{Leia atentamente cada questão!} As questões objetivas terão uma,
                    e apenas uma única, resposta a ser marcada.
                \item \textbf{Desligue o celular} antes de começar. A avaliação é
                    \textbf{individual} e \textbf{sem consulta}.
                \item As questões podem ser respondidas, na prova, com lápis ou caneta. Ao final
                    da prova \textbf{{VOCÊ DEVERÁ PREENCHER O GABARITO COM CANETA
                    DE COR PRETA}} \textbf{\underline{OU AZUL}} para que seja feita a correção
                    automatizada através da leitura do padrão de respostas marcado no
                    gabarito.
                \item \textbf{CUIDADO ao preencher o gabarito} pois eles são \textbf{INDIVIDUAIS
                    e NOMEADOS} (cada estudante tem um gabarito próprio específico, com um
                    código de identificação único). \textbf{Questões rasuradas são anuladas}
                    pelo software de leitura do gabarito.
                \item Ao final da prova, \textbf{DEVOLVA A PROVA E O GABARITO} para o professor.
                \item Siga todas as normas de \textbf{Integridade Acadêmica} da disciplina.
                    Alunos que forem flagrados com qualquer espécie de ``cola'' ou trocando
                    informações com outros alunos terão suas avaliações recolhidas, as notas
                    zeradas, e a situação será encaminhada para a coordenação para a aplicação
                    das penalidades previstas pela universidade.
                \item Com 90 minutos de avaliação você terá tempo de até 8,min por questão,
                    em média.
                \item Boa avaliação!
            \end{itemize}
            \vspace{2cm}
        }}}
    \end{center}
    \newpage

    %%%%%%%%%%%%%%%%%%%%%%%%%%%%%%%%%%%%%%%%%%%%%%%%%%%%%%%%%%%%%%%%%%%%%%%%%%%%%%%%
    %%% Inicia o bloco de questões:
    \begin{questions}

    \newpage
    \question

Em relação aos metadados de um banco de dados, assinale a afirmação correta:
\begin{choices}
\choice Como são sempre atualizados de forma automática pelo sistema de gerenciamento de bancos de dados, estão sempre prontos para nos dar uma ``foto' da estrutura do banco de dados em um momento no tempo.
\choice São armazenados no catálogo (ou dicionário de dados) do sistema de gerenciamento de bancos de dados, que fica nas mesmas tabelas de dados dos usuários, permitindo assim que os metadados estejam sempre atualizados.
\choice São atualizados pelo SGBD, sempre que ocorre alguma operação de inserção, atualização ou remoção de dados.
\choice Armazenam informações sobre os dados dos usuários, ou seja, são dados de rotina dos usuários, aqueles dados que os usuários precisam para trabalhar no seu dia a dia.
\choice Armazenam informações como os dados de cadastro dos usuários, os valores de compra em uma ordem de compra, as turmas e disciplinas que os professores dão aula, e dados de recursos humanos.
\end{choices}

\question

Considere as seguintes tabelas em uma base de dados relacional:\\

\textbf{Departamento (CodDepto, NomeDepto)}\\

\textbf{Empregado (CodEmp, NomeEmp, CodDepto)}\\

Deseja-se obter uma tabela na qual cada linha é a concatenação de uma linha da tabela Departamento com uma linha da tabela de Empregado. Caso um departamento não possua empregados, sua linha no resultado deve conter vazio (NULL) nos campos referentes ao empregado. A operação de álgebra relacional que deve ser aplicada para combinar estas duas tabelas é:
\begin{choices}
\choice Junção interna
\choice Projeção
\choice União
\choice Junção externa
\choice Divisão
\end{choices}

\question

Um banco de dados OLAP (Online Analytical Processing --- Processamento Analítico

Online) é um banco de dados organizado de forma a dar suporte ao Business

Intelligence (BI --- Inteligência de Negócio) que, de forma simplificada, é o

processo de analisar dados para obter informações que podem ser utilizadas na

tomada de decisões comerciais e estratégicas. Em relação aos bancos de dados

OLAP, marque a afirmação correta:
\begin{choices}
\choice São preparados para responder perguntas do tipo ``Quais alunos estão matriculados na turma de Bancos de Dados I neste semestre?'
\choice Geralmente os dados de origem do OLTP são de bancos de dados OLAP.
\choice São otimizados para o processamento de grande volume de transações diárias, o que permite a análise e a tomada de decisões em tempo real, favorecendo assim a administração empresarial.
\choice São otimizados para o processamento de grande volume de dados históricos e métricas de negócio para a tomada de decisões estratégicas, permitindo e facilitando a análise de dados para a geração de informações e conhecimentos para a tomada de decisões.
\choice São otimizados para consulta e relatórios de rotina de dados.
\end{choices}

\question

(Adaptado do ENADE 2005) Avalie o seguite diagrama relacional de um banco de dados:

\begin{figure}[H]

  \centering

  \caption{Banco de dados de peças e fornecedores}

  \label{fig:pecas}

  %\fbox{

    \includegraphics[width=0.5\linewidth]{static/imgs/defaul.png}

  %}\\

  %\footnotesize{Fonte: xxxx.}

\end{figure}

Assinale a opção que apresenta um comando SQL que permite obter uma lista que

contenha o nome de cada fornecedor que tenha fornecido alguma peça, o código da

peça fornecida, a descrição dessa peça e a quantidade fornecida da referida

peça.
\begin{choices}
\choice \begin{verbatim}SELECT f2.nome, p.codigo, p.descricao, f.quantidade FROM pecas INNER JOIN fornecimentos f ON (p.codigo = f.cod_peca) INNER JOIN fornecedores f2 ON (f.cod_forn = f2.cod_forn);\end{verbatim} 
\choice \begin{verbatim}SELECT nome, codigo, descricao, quantidade FROM pecas INNER JOIN FORNECEDORES INNER JOIN FORNECIMENTOS;\end{verbatim} 
\choice \begin{verbatim}SELECT * FROM pecas INNER JOIN fornecedores INNER JOIN fornecimentos;\end{verbatim} 
\choice \begin{verbatim}SELECT * FROM pecas INNER JOIN fornecimentos f ON (p.codigo = f.cod_peca) INNER JOIN fornecedores f2 ON (f.cod_forn = f2.cod_forn) ORDER BY f2.nome;\end{verbatim} 
\choice \begin{verbatim}SELECT DISTINCT p.nome, p.codigo, p.descricao, f.quantidade FROM pecas p INNER JOIN fornecimentos f ON (p.codigo = f.cod_peca);\end{verbatim} 
\end{choices}

\question

Avalie o banco de dados de professores e alunos, mostrado a seguir. Esse banco

de dados pretende armazenar todas as turmas, sendo que um professor pode dar

aula para vários alunos ao mesmo tempo, e um mesmo aluno pode ter aula com

diferentes professores ao mesmo tempo.

\begin{figure}[H]

  \centering

  \caption{Professores, turmas e alunos}

  \label{fig:professores}

  %\fbox{

    \includegraphics[width=0.5\linewidth]{static/imgs/defaul.png}

  %}\\

  %\footnotesize{Fonte: xxxx.}

\end{figure}

Esse banco de dados resolveráo problema de armazenar os dados das turmas?
\begin{choices}
\choice Sim, pois o diagrama mostra claramente que um professor pode ter zero ou mais turmas, e cada aluno pode participar de zero ou mais turmas também.
\choice Não, pois o relacionamento entre 'turmas' e 'professores' está errado: a PK 'cod\_turma' deveria ter ido para a tabela 'professores' como uma FK, mas no diagrama está ao contrário.
\choice Não é possível determinar somente avaliando o diagrama exibido, seriam necessárias mais informações do dicionário de dados.
\choice Sim, pois as cardinalidades e os sentidos dos relacionamentos entre 'professores' e 'alunos' indica, no final, um relacionamento N:N entre essas duas tabelas.
\choice Não, pois as cardinalidades e os sentidos dos relacionamentos entre 'professores' e 'turmas', e entre 'turmas' e 'alunos' não cria um relacionamento N:N entre 'professores' e 'alunos'.
\end{choices}

\question

(Adaptado do ENADE 2005) Analise o relacionamento existente entre as tabelas

``clubes' e ``jogadores', nas alternativas abaixo, e marque a alternativa que

indica que todo jogador deve pertencer a um único clube, e todo clube poder ter

ou não jogadores.
\begin{choices}
\choice \includegraphics[width=0.5\linewidth]{static/imgs/defaul.png} \vspace{0.5cm}
\choice \includegraphics[width=0.5\linewidth]{static/imgs/defaul.png} \vspace{0.5cm}
\choice \includegraphics[width=0.5\linewidth]{static/imgs/defaul.png} \vspace{0.5cm}
\choice \includegraphics[width=0.5\linewidth]{static/imgs/defaul.png} \vspace{0.5cm}
\choice \includegraphics[width=0.5\linewidth]{static/imgs/defaul.png} \vspace{0.5cm}
\end{choices}

\question

Uma definição lógica, abstrata e auto-contida dos objetos que modelam a

estrutura de armazenamento de dados, dos operadores que modelam o comportamento

e dos demais conceitos que modelam a integridade, e que, juntos, constituem a

máquina abstrata com a qual os usuários interagem é chamada de:


\begin{choices}
\choice Modelo entidade-relacionamento
\choice Modelo de dados
\choice Modelo relacional
\choice Nenhuma das respostas acima
\choice Banco de dados
\end{choices}

\question

Quando Edgar Codd cricou o modelo relacional de dados,

também criou algumas regras para que a estrutura das tabelas fosse adequada às

operações de inserção, atualização, remoção e busca de dados. Ao avaliar um

banco de dados que armazena dados dos eleitores, dos partidos e dos candidatos

de uma eleição, você encontrou a situação ilustrada na tabela abaixo:

\begin{figure}[H]

  \centering

  \includegraphics[width=0.5\linewidth]{static/imgs/defaul.png}

\end{figure}

Em relação às boas normas de projeto do modelo relacional, assinale a

alternativa que corretamente identifica um problema no projeto desta tabela:


\begin{choices}
\choice Apesar do telefone ser um atributo multivalorado, do jeito que a tabela está criada só é possível armazenar dois números de telefones, criando uma restrição ao armazenamento de telefones de contado de eleitores que tenham mais de dois telefones.
\choice Todas as afirmativas anteriores.
\choice A tabela contém um atributo multivalorado (o correto seria que cada tabela não tivesse atributos multivalorados).
\choice A tabela está misturando dados de três entidades diferentes, eleitores, partidos e candidatos (o correto seria que cada tabela armazenasse dados de um único assunto ou entidade).
\choice A chave primária para a tabela não está definida e, somente olhando as colunas, não é possível dizer se é a coluna 'numero' ou a coluna 'titulo'.
\end{choices}

\question

Dentre as definições a seguir, ligadas ao conceito de normalização do modelo relacional, qual delas é \textbf{INCORRETA}?


\begin{choices}
\choice O objetivo da normalização é eliminar várias anomalias (ou aspectos indesejáveis) de uma relação.
\choice As formas normais se baseiam em certas estruturas de dependências.
\choice A primeira forma normal estabelece que os atributos da relação contêm apenas valores atômicos.
\choice As relações que obedecem à primeira forma normal não apresentam anomalias.
\choice A normalização é um processo passo a passo reversível de substituição de uma dada coleção de relações por sucessivas coleções de relações as quais possuem uma estrutura progressivamente mais simples e mais regular.
\end{choices}

\question

Um sistema de gerenciamento de bancos de dados (SGBD) é um software servidor

que nos permite definir, construir, manipular, integrar e compartilhar bancos

de dados entre diversos usuários e aplicações. São exemplos de

SGBD, \textbf{exceto}:
\begin{choices}
\choice SQL Server
\choice Oracle Database
\choice PostgreSQL
\choice MongoDB
\choice Oracle SQL Developer
\end{choices}

\question

Na linguagem SQL, qual das seguintes palavras-chave é usada para modificar os

dados em uma tabela existente?


\begin{choices}
\choice SELECT
\choice Nenhuma das alternativas acima
\choice INSERT
\choice CREATE
\choice UPDATE
\end{choices}

\question

Ao executar o seguinte comando SQL, nenhum dado foi retornado, ou seja, não

ocorreu nenhum erro de sintaxe mas nenhum dado foi encontrado:

\begin{tcolorbox}

 \begin{verbatim}1   SELECT c.nome AS cliente

 2   FROM clientes c

 3   LEFT JOIN pedidos p ON (p.cliente_id = c.cliente_id)

 4   WHERE p.cliente_id IS NULL

 5   ORDER BY cliente;

 \end{verbatim}

\end{tcolorbox}

Por que esse comando SQL não retornou nenhum dado?


\begin{choices}
\choice A junção que deveria ter sido realizada, na linha 3, era uma RIGHT JOIN ao invés de uma LEFT JOIN.
\choice Há um erro semântico na linha 3
\choice Há um erro lógico na linha 4.
\choice A junção que deveria ter sido realizada, na linha 3, era uma FULL JOIN ao invés de uma RIGHT JOIN.
\choice Todos os clientes cadastrados fizeram, pelo menos, um pedido de compra.
\end{choices}



    %%%%%%%%%%%%%%%%%%%%%%%%%%%%%%%%%%%%%%%%%%%%%%%%%%%%%%%%%%%%%%%%%%%%%%%%%%%%%%%%
    %%% Finaliza o bloco de questões:
    \end{questions}
    \end{document}
    